% =====================================================================
% VERSION ENSAE, REDUCTION SFE (phase 3, 1er octobre 2026) : annexe B, parametres et notations.
% Copie a texte constant de chapitres/17b_table_parametres.tex puis chapitres/16_notations.tex, les deux anciennes annexes devenant les sections B.1 et B.2.
% Ne changent que les titres et les renvois repares, listes dans
% reduction/phase3.md.
% Lu par main_ensae.tex seul : ne jamais le faire lire par main_institut.tex,
% qui garde chapitres/.
% =====================================================================

% =====================================================================
% TABLE DES PARAMÈTRES  --  annexe.
%
% REMISE DANS LE MÉMOIRE LE 18 SEPTEMBRE 2026, quelques heures après le déport
% des annexes B, D et E. Elle vivait au \S sec:ann-parametres de l'annexe D
% (17_pieces_justificatives.tex) et elle est partie avec elle ; Kélian a demandé
% de la remettre le jour même. Elle a été RETIRÉE du fichier déporté, pas
% recopiée : le même tableau à deux endroits divergerait.
%
% POURQUOI ELLE SEULE REVIENT, quand les vingt et une autres sections de
% l'annexe D restent dehors. C'est la table qu'un jury lit en premier pour
% savoir ce qui est calibré, ce qui est posé et ce qui n'est pas identifié, et
% le chapitre 11 annonce qu'elle existe « pour qu'aucun chiffre du mémoire ne
% soit sans origine déclarée » : la laisser dehors rendait cette phrase fausse.
%
% ELLE EST EN longtable ET NON EN tabular, depuis le 17 août 2026. Un tabular ne
% peut pas se couper entre deux pages : il débordait de 102 pt et hyperref
% émettait l'annotation hors page qui accompagne toujours ce débordement.
%
% SES NEUF SCRIPTS SOURCES SONT NÉCESSAIRES. La ligne
% « % SOURCES-SCRIPTS: 08b 22 25 34 47 60 67 99 43 » couvre les soixante et
% quelques nombres de la table ; en retirer un sort des valeurs du contrôle.
% =====================================================================

\chapter{Paramètres et notations}
\label{sec:ann-parametres}

\section{Nature et sensibilité des paramètres}

Trois natures sont distinguées : calibré sur données (avec son incertitude),
\emph{posé} par jugement d'expert ou choix normatif (non identifiable, cf.\
chapitre~\ref{ch:identifiabilite}), et semi-ancré sur un proxy externe non
spécifique au périmètre.

\vspace{0.25cm}
% PASSEE DE tabular A longtable LE 17 AOUT 2026. Cette table tenait sur une page tant que rien
% ne la precedait dans le chapitre ; l'ajout de la section sur le dispositif de verification l'a
% decalee et elle a debordé de 102 pt, avec l'annotation hyperref hors page qui accompagne
% toujours ce debordement. Un tabular NE PEUT PAS se couper entre deux pages : le piege est
% documente et il vient de mordre. longtable le peut, et l'en-tete se repete.
{\footnotesize
\renewcommand{\arraystretch}{1.28}
\begin{longtable}{@{}>{\raggedright\arraybackslash}p{2.7cm}
                    >{\raggedright\arraybackslash}p{2.6cm}
                    >{\raggedright\arraybackslash}p{1.7cm}
                    >{\raggedright\arraybackslash}p{4.0cm}
                    >{\raggedright\arraybackslash}p{3.6cm}@{}}
\toprule
\textbf{Paramètre} & \textbf{Valeur} & \textbf{Nature} & \textbf{Justification} & \textbf{Sensibilité du besoin} \\
\midrule
\endfirsthead
\toprule
\textbf{Paramètre} & \textbf{Valeur} & \textbf{Nature} & \textbf{Justification} & \textbf{Sensibilité du besoin} \\
\midrule
\endhead
\bottomrule
\endlastfoot
Queue de sévérité $\xi$ & $0{,}5954$ (OpRisk, résultat principal) ; PRC $1{,}03$ ; $0{,}90$ posé en pire cas de modèle & calibré, incertain & GPD sur pertes réelles ; hétérogène par catégorie ($0{,}87$--$1{,}37$) & \textbf{forte} : de $0{,}60$ à $0{,}38$ quand le seuil monte au-dessus de celui retenu, dans l'IC90 $[0{,}30\,;0{,}83]$ ; pilote la queue \\
Échelle GPD $\sigma$, seuil $u$, $p_u$ & $57{,}97$ M\euro{} ; $20{,}03$ M\euro{} ; $0{,}15$ & calibré & POT au percentile $85$ de l'échantillon 2000--2026, $86$ excès sur $570$ pertes ; $p_u=86/570$ (\S\ref{sec:pu-coherence}) & modérée (via la queue) \\
Fréquence $\lambda$ & ${\sim}21{,}6$/an (OpRisk) & convention de calage & $582$ événements collectés / 27 ans, dont douze d'avant 2000 : $+2{,}1\,\%$ sur l'échantillon (\S\ref{sec:pu-coherence}) & sous-linéaire sur le niveau \\
Multiplicateur de fréquence par état & $\times 1{,}525$ partiel ; $\times 2{,}487$ non conforme ($\lambda$ de $21{,}56$ à $53{,}62$) & semi-ancré & parts Hackmageddon et fourchettes posées sur des effets publiés (chapitre~\ref{chap:multietats}) ; aux bornes, $\lambda$ de $41{,}7$ à $65{,}6$ & \textbf{forte} : canal dominant du résultat principal \\
Sur-dispersion $\phi$ & $9{,}20$ & \textbf{posé} (prudent) & $2{,}24$ mesuré sur la série annuelle détendancée, maintenu entre scénarios & élasticité $-0{,}44$ sur l'écart entre états \\
Choc commun fréquence $a$ & $0{,}60$ & \textbf{posé} (prudent) & dépasse l'IC empirique haut ($0{,}22$) : conservateur & enveloppe le modèle à paquets ($\times 2{,}5$) \\
Structure $W = g\,T/s_{\max}$ & classeur qualitatif & \textbf{posé} / normatif & jugement d'expert consigné ; direction non identifiable & l'asymétrie départage les deux piliers de tête ; l'essentiel du surcoût de contagion vient de la partie symétrique \\
Gain de propagation $g$ & $0{,}90$ base ; par état $0{,}45/0{,}68/0{,}90$ & \textbf{posé} & amplitude de $W$, bornée par Leontief $\rho(W)\le g<1$ & \textbf{levier} : $\times 1{,}6$ sur $[0{,}05\,;0{,}90]$ \\
Score de conformité $q$ / états & continu $[0,1]$ (ou NC/PC/C) & \textbf{posé} & décrit le profil, subsume 3 ou 5 niveaux & \textbf{levier} : $-53\,\%$ de NC à C sur le seul score ; facteur $3{,}34$ sur les quatre canaux \\
Multiplicateur de matérialité par état & $0{,}85/1{,}00/1{,}20$ & \textbf{posé} & module $p_u$ selon la conformité & modérée (canal détection) \\
Charge systémique $\gamma$ & $0{,}68$ & semi-ancré & proxy concentration cloud (AWS+Azure+GCP, Synergy 2025) & modérée \\
Seuil de crise $\theta_{\text{crise}}$ & $-2{,}5$ & \textbf{posé} & déclenche le régime de co-mouvement extrême & agit sur la seule queue crise \\
Probabilités d'état $P(\text{NC/PC/C})$ & $0{,}35/0{,}35/0{,}30$ & semi-ancré & enquêtes conformité DORA ${\sim}50\,\%$ non conformes (Deloitte, ESA) & déplace le mélange, pas les états \\
Copule de dépendance $\theta$ & Gumbel $1{,}8$ & \textbf{posé} & queue supérieure entre briques & \textbf{négligeable} : $\Delta_{\text{DORA}}$ $<1{,}4\,\%$ sur $[1,4]$ \\
Excès de source $u_{ij}$ & $\operatorname{logit}p_{ij}=\operatorname{logit}p_j+u_{ij}$ & \textbf{non identifié} & la direction elle-même ; normalisation $\sum_i u_{ij}=0$ (repérage, non information) & mesurée par l'échelle de repli (ch.~\ref{chap:partielle}) : rang~$1$ capte $74\,\%$ et resserre les bornes de $25\,\%$ ; à $u=0$ P1 et P4 deviennent indiscernables \\
Élasticité sévérité/taille $b$ & $0{,}087$ $[0{,}026\,;0{,}148]$ & calibré, tronqué & EMV lognormale tronquée au seuil de collecte (ch.~\ref{chap:donnees}) & faible : $\times 0{,}86$ sur la sévérité transposée \\
Plafond de sévérité & $40$ M\euro{} & \textbf{posé} & capacité de réassurance ($\xi\ge 1$) ; appliqué par pilier touché & borne la queue extrême \\
\end{longtable}
}

\vspace{4pt}
% SOURCES-SCRIPTS: 08b 22 25 34 47 60 67 99 43 76

Les paramètres à forte sensibilité sont soit calibrés sur données (la queue $\xi$, dont
l'incertitude est chiffrée), soit des \emph{leviers voulus} ($g$, $q$, le multiplicateur de
fréquence), qui portent l'effet de la conformité et de la contagion sur le capital. Les paramètres purement posés qui
ne sont pas des leviers ($\theta$ de copule, $\gamma$, $\theta_{\text{crise}}$, probabilités d'état,
matérialité) ont une influence faible à modérée dans ces tests. Les correspondances entre conformité
et canaux restent posées. Leur effet sur l'écart est déterminant ; l'unité de perte et le bootstrap
conjoint en précisent la portée (\S\ref{sec:institut-unite-incertitude}).

\section{Table des notations}

% MIGRÉ depuis memoire/main.tex (plages 3230-3259).
% RELECTURE FAITE (août 2026) : la table héritée ne portait AUCUN symbole du
% modèle de cascade (W, S, A, TRANS, ROOT, g, s_j, e_j, R_0, t, kappa), alors
% que c'est le modèle central. Elle est réécrite et organisée en cinq blocs.
\label{not:ann:notations}

\noindent Les symboles sont regroupés par brique du modèle. Les probabilités, scores et gains sont
sans unité ; les pertes sont en M\euro{} et les fréquences en incidents par an. Les indices
$i$, $j,k$, $n$ et $y$ désignent respectivement l'entité, les piliers, l'incident et l'année.
Certains symboles ont un sens local : $A$ désigne une matrice antisymétrique ou des actifs,
$S$ une matrice symétrique ou une coalition, et $\rho$ une corrélation ou un rayon spectral.

% LONGTABLE ET NON TABULAR. Un tabular ne peut pas se couper : la table dépassait le
% bas de page, et hyperref plaçait alors les liens \ref des sept dernières lignes à des
% coordonnées NÉGATIVES (sept avertissements « Annotation out of page boundary » à la
% compilation). Les renvois fonctionnaient mais pointaient sur la mauvaise page du PDF.
% longtable répartit la table sur plusieurs pages et rejoue l'en-tête à chaque page.
\renewcommand{\arraystretch}{1.12}
{\small
\begin{longtable}{>{\raggedright}p{3.3cm} >{\raggedright\arraybackslash}p{10.2cm} c}
\toprule
\rowcolor{lightblue}
\textbf{Symbole} & \textbf{Signification} & \textbf{Ch.} \\
\midrule
\endfirsthead
\toprule
\rowcolor{lightblue}
\textbf{Symbole} & \textbf{Signification} & \textbf{Ch.} \\
\midrule
\endhead
\bottomrule
\multicolumn{3}{r@{}}{\footnotesize\itshape suite page suivante} \\
\endfoot
\bottomrule
\endlastfoot
\multicolumn{3}{@{}l}{\emph{\textcolor{navy}{Piliers et cascade dirigée}}} \\
$P_1,\dots,P_5$ & Les cinq piliers DORA : gouvernance, gestion des incidents, tests de résilience, prestataires tiers, partage d'informations & \ref{chap:cascade} \\
$T=(t_{jk})$ & Matrice d'influence brute du classeur qualitatif, asymétrique, non normalisée & \ref{chap:cascade} \\
$\mathrm{ROOT}$ & Propension de chaque pilier à \emph{amorcer} une chaîne ; répartit les amorces du moteur & \ref{chap:cascade} \\
$\mathrm{GBASE}$ & Gravité intrinsèque de chaque pilier s'il cède seul, barème ordinal du modèle qualitatif, absent du moteur de capital & \ref{chap:cascade} \\
$s_j$ & Somme de la ligne $j$ de $T$ ; $s_{\max}=\max_j s_j$ & \ref{chap:cascade} \\
$g$ & Gain de propagation, somme maximale des probabilités émises par un pilier dans la matrice normalisée de référence ; sans unité & \ref{chap:cascade} \\
$W$ & Matrice de contagion normalisée à la Leontief, $W=g\,T/s_{\max}$. \textbf{Convention : $W_{jk}$ va de la source $j$ vers la cible $k$ ; la ligne émet, la colonne reçoit} & \ref{chap:cascade} \\
$e_j$ & Propension à propager depuis $j$, $e_j=g\,s_j/s_{\max}$ & \ref{chap:cascade} \\
$\rho(W)$ & Rayon spectral de $W$ ; $\rho(W)\le g<1$ garantit la stabilité & \ref{chap:cascade} \\
$(I-W)^{-1}$ & Inverse de Leontief, somme des marches pondérées avec multiplicité ; $I$ est l'identité, $W^\top$ la transposée & \ref{chap:cascade} \\
$M,\ R_0$ & Matrice de génération suivante et son rayon spectral (extinction si $R_0<1$) & \ref{chap:cascade} \\
$S_a$ & Ensemble des piliers atteints depuis l'amorce $a$ & \ref{chap:cascade} \\
$K_j,\ u_j$ & Seuil latent d'incident ; seuil monétaire POT. Ils diffèrent des critères réglementaires DORA (art.~18) & \ref{chap:cascade} \\
$X_i,\ X_{ij}$ & Stress latent de l'entité $i$, puis de son pilier $j$ ; sans unité, distinct de la perte monétaire $X_n$ & \ref{chap:cascade} \\
$\Phi,\ \Phi^{-1}$ & Répartition et fonction quantile de la normale standard ; $\mathbf 1\{E\}$ vaut un si $E$ se réalise & \ref{chap:cascade} \\
\midrule
\multicolumn{3}{@{}l}{\emph{\textcolor{navy}{Identifiabilité et identification partielle}}} \\
$S,\ A$ & Parties symétrique et antisymétrique de $W$ ; $S$ est posé à dire d'expert, $A$ est exploré sous contraintes & \ref{ch:identifiabilite} \\
$J_{jk}$ & Courant de probabilité $\pi_jP_{jk}-\pi_kP_{kj}$ ; nul ssi réversibilité & \ref{ch:identifiabilite} \\
$\pi,\ \tilde P$ & Loi stationnaire ; noyau adjoint (chaîne renversée dans le temps) & \ref{ch:identifiabilite} \\
$\mathcal{E}(f,f)$ & Forme de Dirichlet de la chaîne auxiliaire stationnaire ; ne dépend que de la symétrisée pondérée $S_\pi$ & \ref{ch:identifiabilite} \\
$\sigma$ & Production d'entropie ; nulle ssi la chaîne est réversible & \ref{ch:identifiabilite} \\
$\Delta(N),\ z$ & Statistique d'asymétrie de la matrice $N$ des transitions observées et son écart réduit au placebo par permutation & \ref{ch:identifiabilite} \\
$t$ & Degré d'ignorance directionnelle, $|A_{jk}|\le t\,S_{jk}$ ; $t=1$ = ignorance totale & \ref{chap:partielle} \\
$\mathcal{W}(t)$ & Ensemble admissible conditionnel au $S$ d'expert et aux contraintes posées ; les extrema aux sommets ne sont pas des bornes globales démontrées & \ref{chap:partielle} \\
$u_{ij}$ & Excès de propagation propre à la source, $\operatorname{logit}p_{ij}=\operatorname{logit}p_j+u_{ij}$ & \ref{chap:partielle} \\
$L(\mathcal{V}),\ V(\mathcal{I})$ & Largeur de bande de capital ; valeur d'une information & \ref{ch:livrable} \\
\midrule
\multicolumn{3}{@{}l}{\emph{\textcolor{navy}{Socle : fréquence et sévérité}}} \\
$\xi,\ \sigma_{\text{GPD}},\ u$ & Indice de queue, échelle et seuil POT de la GPD de sévérité & \ref{chap:socle} \\
$\zeta_u=p_u$ & Probabilité de dépassement du seuil par une sévérité élémentaire, $\Prob(X>u)$ & \ref{chap:socle} \\
$N_u$ & Nombre d'excès au-delà de $u$ & \ref{chap:socle} \\
$\lambda_{ref}$ & Fréquence de référence PRC ($341$/an) : \emph{clé de répartition} par vecteur, n'alimente pas le moteur & \ref{chap:donnees} \\
$\lambda$ & Fréquence d'entrée, événements TIC enregistrés ($21{,}6$/an au secteur), avant le tirage de la perte par pilier & \ref{chap:donnees} \\
$\phi\ (\varphi),\ a$ & Rapport variance/moyenne du comptage ; charge commune du mélange lognormal & \ref{chap:socle} \\
$r,\ p$ & Forme et probabilité de la binomiale négative ($p$ distinct de $p_u$) & \ref{chap:socle} \\
$N_y,\ Y_{n,k},\ X_n,\ L_y$ & Incidents de l'année ; perte du pilier $k$ dans l'incident $n$ ; perte de l'incident ; somme annuelle & \ref{chap:socle} \\
$b$ & Élasticité de la sévérité à la taille de la firme & \ref{chap:donnees} \\
\midrule
\multicolumn{3}{@{}l}{\emph{\textcolor{navy}{Conformité multi-états}}} \\
$C^\star_j,\ \Theta,\ \varepsilon_j$ & Capacité latente du pilier $j$ ; facteur commun ; choc propre (normales réduites indépendantes) & \ref{chap:multietats} \\
$U_j$ & Score de conformité $\Phi(C^\star_j)\in[0,1]$, de coupures $0{,}35$ et $0{,}70$ & \ref{chap:multietats} \\
$\gamma$ & Charge systémique de la latente de conformité ($0{,}68$) & \ref{chap:multietats} \\
$\omega$ & Probabilité du choc commun d'accumulation ($0$ en C, $0{,}68$ en NC) & \ref{chap:multietats} \\
$K_{\text{bas}},K_{\text{haut}}$ & Seuils ordonnés découpant les états NC / PC / C & \ref{chap:multietats} \\
$Q$ & Générateur de la chaîne de mise en conformité (séjours de type phase) & \ref{chap:multietats} \\
$v_0,\ \overline B(t)$ & Probabilités initiales des phases ; capital moyen conditionnel en $t$, distinct du quantile du mélange & \ref{chap:multietats} \\
$S_0,S_1,S_2$ & Scénarios C / PC / NC (multiplicateurs de fréquence) & \ref{chap:socle} \\
\midrule
\multicolumn{3}{@{}l}{\emph{\textcolor{navy}{Capital, allocation et décision}}} \\
$\VaR_\alpha(L)$, $\TVaR_\alpha(L)$ & Value-at-Risk et Tail-VaR de la charge \textbf{annuelle agrégée} $L$, au niveau $\alpha=99{,}5\,\%$ : quantile de capital et diagnostic de moyenne de queue & \ref{chap:socle} \\
$B,\ B_e,\ B(W),\ B(A)$ & Besoin ORSA du modèle, $\VaR_{99{,}5\%}(L)$, sous un état $e$, une matrice $W$ ou à des actifs $A$ ; pas un SCR réglementaire & \ref{chap:multietats} \\
$\qsev_\alpha(X),\ \qbarsev_\alpha(X)$ & Quantile d'un sinistre unique et moyenne des pertes au-delà ; grandeurs descriptives distinctes du capital annuel & \ref{chap:socle} \\
$\Delta_{DORA}$ & Écart de capital entre états NC et C supposés, sans lecture causale & \ref{chap:resultats} \\
$C_j,\ AC_j,\ \varphi_j$ & Contributions à la TVaR en queue ; Shapley du pilier $j$ (distinct de $\phi$) & \ref{chap:resultats} \\
$A,\ A_m,\ \Lambda$ & Actifs de l'entité ; actifs médians des pertes ; incidents annuels moyens du panel & \ref{chap:resultats} \\
$b_\lambda,\ m(A)$ & Exposant de fréquence aux actifs ; facteur de sévérité à la taille de l'entité & \ref{chap:resultats} \\
$D,\ d,\ \tau,\ \mathrm{Adj}$ & Écart ajouté ; $D/\mathrm{BSCR}$ ; part retenue en agrégation corrélée ; absorption des pertes & \ref{chap:resultats} \\
$\mathrm{LR}$ & Ratio de sinistralité, pertes attendues sur primes (sans unité) & \ref{chap:resultats} \\
$\mathrm{CoC}$ & Coût du capital de la marge de risque ($6\,\%$, puis $4{,}75\,\%$) & \ref{chap:resultats} \\
$\ell,\ \ell^\star$ & Portée d'un traité en excès de perte annuel ; portée optimale & \ref{chap:resultats} \\
$\kappa$ & Coefficient d'inefficacité du transfert (capacité, base, contrepartie) & \ref{chap:resultats} \\
$\theta$ & Paramètre de la copule de Gumbel (comparaison, non modèle) & \ref{ch:robustesse} \\
\end{longtable}
}

\clearpage

% =====================================================================

% SOURCES-SCRIPTS: 63 67 22 50
